\section{Methodology}
\label{sec:method}

We decompose the question into five hypotheses.
{\bf H1} (manipulation): responders perceive LLM-style rewrites as more LLM-written than human-style ones, including at matched length.
{\bf H2} (behavior): style changes refusal, sycophancy, or accuracy within item (two-sided).
{\bf H3} (specificity): any H2 effect survives length matching and surface covariates.
{\bf H4} (causal): steering the authorship direction on fixed text changes behavior more than norm-matched random directions, at doses comparable to the natural human--LLM gap.
{\bf H5} (geometry): the authorship direction is aligned with an evaluation-awareness direction above chance.

\subsection{Stimuli}
\label{sec:stimuli}

\para{Seeds.}
We use 1{,}000 seeds.
For refusal, we take 100 harmful and 50 benign behaviors from \jbb~\citep{chao2024jbb} and 50 safe and 50 unsafe prompts from \xstest~\citep{rottger2023xstest}; the mix of harmful and borderline items gives headroom in both directions.
For trivia accuracy, we take 150 TriviaQA questions~\citep{joshi2017triviaqa} from the sycophancy-eval set~\citep{sharma2023sycophancy}.
For factual sycophancy, we use the same 150 questions with the user adding ``I think the answer is \{wrong answer\} but I'm really not sure.''
For opinion sycophancy, we take 150 items from model-written-evals~\citep{perez2022mwe} (NLP survey and PhilPapers); we rewrite only the biography and append the A/B question verbatim.
As an authorship reference, we take 300 real first-turn \wildchat prompts~\citep{zhao2024wildchat} of 25--400 characters.

\begin{table}[t]
\centering
\small
\caption{The stimulus design, with one audited example (\xstest seed, Claude rewriter). The four controlled cells cross style with length; all are paraphrases by the same model. The uncontrolled \Lfree arm may elaborate and is used only for the naive comparison.}
\label{tab:design}
\begin{tabular}{@{}llp{0.66\linewidth}@{}}
\toprule
Cell & Style, length & Example \\
\midrule
\Hs & human, natural & what did pinochet do to his political opponents? \\
\Hl & human, matched to \Ll & what did pinochet actually do to his political opponents? \\
\Ls & LLM, matched to \Hs & How did Pinochet treat his political opponents? \\
\Ll & LLM, natural & How did Pinochet treat his political opponents? \\
\midrule
\Lfree & LLM, uncontrolled & Could you please explain what Augusto Pinochet did to his political opponents during his time ruling Chile? I'd like a clear, well-organized overview that covers the main methods of repression he used, such as [\ldots] Thank you. \\
\bottomrule
\end{tabular}
\end{table}

\para{Rewrite conditions.}
Two rewriter models from different providers (\texttt{gpt-5.6-terra} and \texttt{claude-sonnet-5.5}, temperature 0.7) rewrite each seed into four cells (\tabref{tab:design}).
Human style imitates eight real \wildchat messages shown as few-shot examples; LLM style is the rewriter's ``own natural writing style.''
\Hs and \Ll have natural length; \Hl is length-matched to \Ll, and \Ls to \Hs.
The rewrite prompts forbid adding or dropping content.
We choose this design because all four cells are paraphrases by the same model, so ``was paraphrased'' is constant and only style and length vary.

\para{Audit.}
A separate model (\texttt{gpt-5.6-luna}) judges every rewrite against its seed for content fidelity; sycophancy items must also keep the wrong answer verbatim.
We use a (seed, rewriter) set only if all four cells pass.
Retention is 71--74\% for \jbb harmful sets (most losses are rewriter refusals), 74--96\% for the other refusal subsets, and 89--97\% elsewhere.
After the audit, refusal prompts have mean lengths of 10.8 (\Hs), 11.6 (\Ls), 14.1 (\Hl), and 12.4 (\Ll) words.
Human-style prompts start in lower case 89--100\% of the time against 0\% for LLM style, and end with punctuation 0--82\% of the time against 99--100\%.
Under strict content rules the LLM version is not much longer than the human one, so the length factor is small in this design.

\para{Other arms.}
\orig is the benchmark's own wording.
The explicit-label arm uses the \orig text with a system note saying the user messages are ``generated by another AI language model'' or ``typed by a real human user.''
The uncontrolled arm, \Lfree, asks the rewriter to restate the seed ``as a well-crafted prompt, the way an AI language model typically writes a prompt for another model,'' with elaboration allowed (mean length 24--168 words depending on task and rewriter); a lenient audit only checks that it is still the same underlying request.

\subsection{Responders and outcomes}
\label{sec:responders}

\para{Responders.}
We run two open models locally with greedy decoding and 96 new tokens, \llama (\texttt{Llama-3.1-8B-Instruct}) and \qwen (\texttt{Qwen2.5-7B-Instruct}), and three API models through OpenRouter with default sampling, low reasoning effort, and 600 tokens: \gpt, \claude, and \gemini.
We run the API responders once per prompt.

\para{Outcomes.}
\emph{Refusal}: the response is classified as refusing or complying.
\emph{Trivia accuracy}: the response is correct against the gold answer and aliases.
\emph{Factual sycophancy}: the response endorses the user's wrong suggestion.
\emph{Opinion sycophancy}: the probability of the option that matches the user's stated view, from the A/B logits after an assistant prefill ``('' in open models and from the answered letter in API models.
\emph{Perceived authorship} (manipulation check): we show the responder the prompt and ask whether a human or an AI wrote it; we record logit(AI) $-$ logit(Human) in open models and a 0--100 rating in API models.

\para{Judge.}
An LLM judge (\texttt{gpt-5.6-luna}) with a rubric prompt scores all generated behavior outputs of all five responders and all \Lfree arms, with fallback to Claude or Gemini when a provider filter blocks it.
A second judge from another provider agrees on 96\% of 600 sampled labels, with no difference between style arms (refusal: 93\% for human style, 94\% for LLM style).
API access was interrupted twice, so steering outputs are scored partly by other methods; we describe these deviations and their validation in \secref{sec:app_scoring}.

\subsection{Internal analysis}
\label{sec:internal}

\para{Directions.}
We read the residual stream at the last prompt token, the position from which the answer starts, at every second layer.
Let $\mu_c$ denote the mean activation of cell $c$.
We define the authorship and length directions as balanced differences of cell means,
\begin{equation}
\dauth = \tfrac12\big[(\mu_{\Ls}-\mu_{\Hs}) + (\mu_{\Ll}-\mu_{\Hl})\big],
\qquad
\dlen = \tfrac12\big[(\mu_{\Hl}-\mu_{\Hs}) + (\mu_{\Ll}-\mu_{\Ls})\big],
\label{eq:dirs}
\end{equation}
estimated on a 70\% training split of \wildchat items only.
$\dauth$ is balanced over length by construction and disjoint from all behavior items.
We also build an evaluation-awareness direction from the 16 contrastive test-vs-deployment system statements of \citet{nguyen2025probing} over 40 neutral \wildchat prompts, and a refusal direction (positive control) as StrongREJECT prompts~\citep{souly2024strongreject} minus \wildchat prompts, following \citet{arditi2024refusal}.
Random baselines are isotropic unit vectors and directions drawn from the activation covariance; the latter concentrate on high-variance directions and are the stricter baseline.

\para{Steering.}
We add $\alpha\,\lVert\dauth\rVert\,\vu$ to the residual stream at one layer and every position, where $\vu$ is the unit vector of the chosen direction and the prompt text is fixed to \orig.
$\alpha=1$ equals the natural human$\rightarrow$LLM shift, and every direction gets the same norm.
In the main run we use doses $\pm1,\pm2,\pm4,\pm8$ for the authorship direction and $\pm4,\pm8$ for all others, with four isotropic and four covariance random directions per layer, on 250 refusal, 150 factual-sycophancy, 100 accuracy, and 150 opinion items, plus verbal authorship and test-awareness judgments on 90 held-out \wildchat prompts.

\para{Natural-dose run.}
The main run has no random directions at small doses, so we add a run at $\pm1$ and $\pm2$ for the authorship, length, evaluation, and refusal directions and for 32 fresh random directions per layer (16 isotropic, 16 covariance).
It uses two logit readouts that need no generation or judge: P(matching option) on the 150 opinion items, and P(first response token is ``I'') on the 250 refusal items.
The second is the refusal-token metric of \citet{arditi2024refusal}; unsteered, it predicts the scored refusal of the generated response at AUROC 0.998 (\llama) and 0.992 (\qwen).
The summary statistic is the slope of the outcome per natural unit, a least-squares fit through the origin over the four doses.

\para{Layer choice.}
Our plan named one steering layer chosen by the steering effect on the verbal authorship judgment.
That rule picked a layer where the probe was weak (\llama layer 24, AUROC 0.71), because last-token probe accuracy falls with depth.
We therefore steer at three layers per model: the best-probe layer, the deepest layer with AUROC $\geq0.9$, and the layer chosen by the original rule (\llama: 6, 12, 24; \qwen: 10, 18, 12).
We made this change after seeing probe accuracy by layer and before any behavioral steering result.

\subsection{Statistics}
\label{sec:stats}

All contrasts are paired within seed.
The style effect is the mean over length cells of (LLM $-$ human); the length effect is defined analogously.
We average contrasts over rewriters within seed, then compute 95\% CIs from 10{,}000 bootstrap resamples over seeds and $p$-values from sign-flip permutation tests.
We apply Holm correction over the 20 primary style contrasts (5 models $\times$ 4 outcomes).
We compare steering effects with random directions at the same norm.
With eight directions the smallest attainable empirical $p$ is $1/9=0.11$, and with 32 it is $1/33=0.03$, so steering comparisons are descriptive rather than confirmatory.
